\section{Problem Formulation and Design Requirements}
\label{sec:problem_formulation}

\subsection{Advisory State and Decision Space}
Let an incoming farmer interaction be denoted $\mathbf{x} = \langle q_t, I_v, \mathcal{M}, \mathcal{H} \rangle$, where $q_t$ is the natural-language query (standard Bengali, a regional dialect, or romanized Banglish), $I_v$ is an optional uploaded crop/leaf photograph, $\mathcal{M}$ is environmental metadata (timestamp, administrative region, season), and $\mathcal{H}$ is the conversation history.

The system maps $\mathbf{x}$ to an action $a(\mathbf{x}) \in \mathcal{A}$:
\begin{equation}
    \mathcal{A} = \left\{ 
    \begin{array}{ll}
        \text{\texttt{CERTIFY}}, & \text{deterministic single-record resolution},\\
        \text{\texttt{GENERATE\_THEN\_VERIFY}}, & \text{grounded generation subject to verification},\\
        \text{\texttt{CLARIFY}}, & \text{conversational disambiguation},\\
        \text{\texttt{ABSTAIN}}, & \text{fail-closed refusal with escalation pointer},\\
        \text{\texttt{ESCALATE}}, & \text{policy block / human (Krishi 16123) handoff}
    \end{array}
    \right\}
\end{equation}

\subsection{The Critical Agricultural Tuple (11-Slot Contract)}
In chemical and biological plant protection, an actionable recommendation is a tuple over an agronomic decision contract $\mathcal{C}$:
\begin{equation}
    \mathcal{C} = \left\langle 
    c, p, s, a, f, d_{\min}, d_{\max}, u_d, v_s, t_{\text{int}}, t_{\text{phi}}
    \right\rangle
\end{equation}
where $c$ is the host crop, $p$ the target pest/pathogen, $s$ the growth stage, $a$ the registered active ingredient, $f$ the formulation, $[d_{\min}, d_{\max}]$ the approved dosage interval, $u$ the dosage unit, $v_s$ the solvent volume denominator (e.g.\ per 10 L water), $t_{\text{int}}$ the spray interval in days, and $t_{\text{phi}}$ the mandatory pre-harvest interval in days. A record may also carry a regulatory polarity flag $\rho \in \{\text{APPROVED}, \text{BANNED/RESTRICTED}\}$; in the certification predicate below we fold $\rho$ into the authority check rather than the tuple, so the tuple remains an 11-slot agronomic contract while the evidence layer governs its authority.

\subsection{Evidence Authority Model}
Let $\mathcal{E}$ be the set of typed evidence records in the fact store. Each record $e \in \mathcal{E}$ carries two logically distinct attributes:
\begin{align}
    \operatorname{Authority}(e) &= g(\text{source},\, \text{jurisdiction}) \nonumber \\
    \operatorname{Validity}(e) &= g(\text{date},\, \text{version},\, \text{regulatory state}) \nonumber
\end{align}
A record is \emph{authoritative} if it originates from an accredited institution (MoA regulatory registry, BARI, BRRI, DAE) with the correct jurisdiction (Bangladesh). A record is \emph{current} if its effective date is not superseded and its regulatory state has not been revoked. Both must hold simultaneously. An old authoritative record can be obsolete; a recent secondary record is not authoritative.

The system certifies the advisory $\mathcal{C}$ only if a single evidence record jointly entails every slot:
\begin{equation}
    \label{eq:certification_rule}
    \operatorname{Certify}(\mathcal{C}) = 1 \iff \exists\, e^* \in \mathcal{E}:\;
    \begin{array}{l}
        \operatorname{Authoritative}(e^*) \;\land\; \operatorname{Current}(e^*)\; \land \\
        \operatorname{JointlyEntails}(e^*, \mathcal{C}) \;\land\; \operatorname{Complete}(\mathcal{C}) \;\land\\
        g(\mathbf{x}) \;\geq\; \theta^*
    \end{array}
\end{equation}
where $\operatorname{JointlyEntails}$ requires all 11 slots to be satisfied by the \emph{same} record (never assembled across records), $\operatorname{Complete}$ requires every safety-critical slot to be non-null, $g(\mathbf{x})$ is the calibrated selective score, and $\theta^*$ is the frozen threshold selected on a development split. If Eq.~\eqref{eq:certification_rule} fails, the system fails closed: it must not complete the recommendation; it clarifies, abstains, or escalates. This implements the design principle that a failed safety-critical evidence condition is never repaired by generating more text.

\subsection{Primary Endpoints}
We use the following primary endpoints throughout, reported with 95\% Wilson confidence intervals where sample sizes permit.
\begin{enumerate}
    \item \textbf{CUAR} --- Critical Unsafe Acceptance Rate: \emph{unsafe cases certified as safe} $/$ \emph{total unsafe evaluated cases}. Primary safety endpoint.
    \item \textbf{CAC} --- Certified Advisory Correctness: \emph{correct certified advisories} $/$ \emph{certified advisories issued}. Primary advisory endpoint.
    \item \textbf{Resolution coverage / advisory coverage / safe certified coverage}: respectively the fraction of queries receiving some output, the fraction of answerable queries receiving useful advice, and the fraction of answerable queries receiving a \emph{correct certified} advisory.
    \item \textbf{AA} --- Appropriate Abstention: \emph{correct abstentions} $/$ \emph{cases requiring abstention}.
    \item Per-slot accuracy for each of the 11 contract fields, reported individually.
\end{enumerate}

\subsection{Reliability Requirements}
Table~\ref{tab:system_requirements} maps each reliability requirement to its failure mode and the fail-closed response required of the architecture. This table is the design contract used in Sections~\ref{sec:system_architecture}--\ref{sec:knowledge_governance} and the evaluation in Section~\ref{sec:experimental_methodology}.

\begin{table*}[pos=t]
\centering
\caption{System reliability requirements, operational failure modes, and bounded-authority containment mechanisms.}
\label{tab:system_requirements}
\small
\resizebox{\textwidth}{!}{\begin{tabular}{llll}
\toprule
\textbf{Reliability Requirement} & \textbf{Potential Failure Mode} & \textbf{Advisory Risk} & \textbf{System Response / Containment} \\
\midrule
Evidence Completeness & Missing PHI / Dosage bounds & Chemical residue / Toxic crop harvest & Fail-closed abstention (\texttt{ABSTAIN}) \\
Relational Integrity & Cross-record pesticide misbinding & Ineffective or phytotoxic application & Single-record joint binding certification \\
Visual Semantic Alignment & Crop/Disease classification error & Misdirected agronomic advice & Metadata routing fallback \& clarification \\
Perception Uncertainty & Low confidence visual input ($\tau < 0.80$) & Erroneous disease identification & Fallback to conversational clarification \\
Temporal Validity & Obsolete chemical recommendation & Application of withdrawn/banned agrochemicals & Temporal invalidation \& policy override \\
Adversarial Robustness & Prompt injection / Delimiter hijack & System instruction override & Deterministic structured extraction filter \\
Channel Bandwidth Limit & SMS truncation / Buffer overflow & Silent loss of vital safety warnings & Deterministic 160-char GSM template \\
Network Degradation & High packet loss / Cloud disconnect & Service unavailability in rural edge & Offline SQLite cryptographic cache \\
\bottomrule
\end{tabular}}
\end{table*}

